% !PRINTED 334
\weberpage{334}
If we stop, to begin with, at the first difference, we obtain as
\gesperrt{first correction}
\weq{7}{ u = \delta\frac{\varDelta}{\varDelta_\alpha}, }
and if we now set
\[ u = \delta\frac{\varDelta}{\varDelta_\alpha} + u' \]
then it follows from (6), if one omits $u'$ in the second term, that
\[ 0 = \frac{\varDelta_\alpha}{\delta}u' +
  \frac{\varDelta'_\alpha \varDelta(\varDelta-\varDelta_\alpha)}{2\varDelta_\alpha^2}, \]
and hence as second correction:
\weq{8}{ u' = \delta\frac{\varDelta'_\alpha
  \varDelta(\varDelta_\alpha-\varDelta)}{2\varDelta_\alpha^3}. }

The first correction is obtained by replacing the arc of the curve
running between $\alpha$ and $\alpha+\delta$ by the chord, as above;
the second by replacing the arc between $\alpha,
\alpha+\delta, \alpha+2\delta$ by a parabola which passes through
the same points, but which now has its \gesperrt{axis parallel to the $x$-axis}.

Let us take as an example the equation
\[ f(x) = x^3 - 2x - 2 = 0, \]
which has a real root between 1,7 and 1,8.

One computes
\[
\begin{array}{llll}
x = 1,7, & f(x) = -0,487, & \varDelta_\alpha = 0,719, & \varDelta'_\alpha = 0,108,\\
\phantom{x = {}}1,8, & \phantom{f(x) {}}{=} +0,232, & \phantom{\varDelta_\alpha {}}{=} 0,827, & \\
\phantom{x = {}}1,9, & & & \phantom{\varDelta'_\alpha {}}{=} 1,059.
\end{array}
\]
Since $f(x) = 0$ is to hold, we must set
\[ \varDelta = 0,487 \]
and the first correction to $\alpha = 1,7$ is, by (7),
\[ u = 0,06773, \]
the second correction gives, by (8),
\[ u' = 0,0164, \]
hence
\[ \alpha + u + u' = 1,76937. \]

The more accurate value, computed in another way, is 1,76929\ldots\ We
thus have a result accurate in the first three decimals. But we have
here no other means of judging the accu-
%CONT>>
